\subsection{关于纯量环有限扩张的性质}

在本节中，A 与 B 表示两个整闭 Noether 整环，满足 \(A \subset B\) ，B 是有限生成 A-模，而 K 与 L 分别是 A 与 B 的分式域。在涉及 A-模的场合，我们将用 \(div_{A}, \chi_{A}, c_{A}, \gamma_{A}, r_{A}\) 分别代替 div, \(\chi, c, \gamma, r\) ，对 B-模采用类似的记号。

我们知道（§ 1，第 10 节）：B 的素理想 P 高为 1，当且仅当 p = P ∩ A 高为 1；此外（出处同上，命题 14），对 p ∈ P(A)，位于 p 之上的素理想 P ∈ P(B) 只有有限多个。为简便起见，我们将用 P\textbar p 表示关系 ``P 位于 p 之上''（即 p = P ∩ A）；随后用 e\_\{P/p\} 或 e(P/p) 表示 v\_\{P\} 在 v\_\{p\} 之上的分歧指数 e(v\_\{P\}/v\_\{p\})（第 VI 章，§ 8，第 1 节），并用 f\_\{P/p\} 或 f(P/p) 表示剩余类次数 f(v\_\{P\}/v\_\{p\})（出处同上）；回忆离散赋值 v\_\{p\} 与 v\_\{P\} 都是赋范的，且 f\_\{P/p\} 是 B/P 的分式域在 A/p 的分式域上的次数。我们置 n = r\_A(B)，其中 B 被视为 A-模；于是按定义 n = {[}L: K{]}，且对一切 p ∈ P(A)，n 也是自由 A\_p-模 B 的秩，对一切 P\textbar p 皆然。于是由第 VI 章，§ 8，第 5 节，定理 2，对一切 p ∈ P(A)：

\[
\sum_ {\mathfrak {P} / \mathfrak {p}} e _ {\mathfrak {P} / \mathfrak {p}} f _ {\mathfrak {P} / \mathfrak {p}} = n.\tag{9}
\]

这样，由于 D(A) 与 D(B) 是自由 Z-模，我们用下列条件定义一个序群的递增同态 N: D(B) → D(A) （也记作 \(N_{B/A}\) ）：

\[
\mathrm{N} (\mathfrak {P}) = f _ {\mathfrak {P} / \mathfrak {p}}. \mathfrak {p} \quad \text {   for   } \quad \mathfrak {P} \in \mathrm{P} (\mathrm{B}), \quad \text {   where   } \quad \mathfrak {p} = \mathfrak {P} \cap \mathrm{A}.\tag{10}
\]

另一方面（§ 1，第 10 节，命题 14），我们已用下列条件定义了序群的递增同态 \(i\colon\mathbf{D}(\mathbf{A})\to\mathbf{D}(\mathbf{B})\) （也记作 \(i_{\mathbf{B}/\mathbf{A}}\) ）：

\begin{enumerate}
\def\labelenumi{(\arabic{enumi})}
\setcounter{enumi}{10}
\tightlist
\item
\end{enumerate}

\[
i (\mathfrak {p}) = \sum_ {\mathfrak {P} / \mathfrak {p}} e _ {\mathfrak {P} / \mathfrak {p}}. \mathfrak {P} \quad \text {   for   } \mathfrak {p} \in \mathrm{P} (\mathrm{A}).
\]

显然，对 A（相应地，B）的除子的每个族 \((d_{\iota})\) （相应地，\((d_{\iota}^{\prime}))\) ）：

\begin{enumerate}
\def\labelenumi{(\arabic{enumi})}
\setcounter{enumi}{11}
\tightlist
\item
\end{enumerate}

\[
\mathrm{N} (\sup (d _ {\iota} ^ {\prime})) = \sup (\mathrm{N} (d _ {\iota} ^ {\prime})), \quad \mathrm{N} (\inf (d _ {\iota} ^ {\prime})) = \inf (\mathrm{N} (d _ {\iota} ^ {\prime}))\tag{13}
\]

\[
i \left(\sup (d _ {\iota})\right) = \sup (i (d _ {\iota})), \quad i (\inf (d _ {\iota})) = \inf (i (d _ {\iota})).
\]

公式 (9) 表明：

\[
\mathbf {N} \circ i = n. 1 _ {\mathbf {D} (\mathbf {A})}.\tag{14}
\]

对一切 \(a \in \mathbf{A}\) （ \(\S 1\) ，第 10 节，命题 14）：

\[
i \left(\operatorname{div} _ {\mathbf {A}} (a)\right) = \operatorname{div} _ {\mathbf {B}} (a).\tag{15}
\]

我们（借助 (13)）推出，对 A 的每个分式理想 a，也有：

\[
i \left(\operatorname{div} _ {\mathbf {A}} (\mathfrak {a})\right) = \operatorname{div} _ {\mathbf {B}} (\mathfrak {a B}).\tag{16}
\]

对每个元素 \(b \in \mathbf{B}\) ，我们知道（第 V 章，§1，第 3 节，命题 11 的推论）\(\mathrm{N}_{\mathrm{L / K}}(b) \in \mathbf{A}\) ；此外（第 VI 章，§8，第 5 节，公式 (9)）：

\[
v _ {\mathfrak {p}} \left(\mathrm{N} _ {\mathrm{L} / \mathbf {K}} (b)\right) = \sum_ {\mathfrak {P} / \mathfrak {p}} f _ {\mathfrak {P} / \mathfrak {p}} v _ {\mathfrak {P}} (b)\tag{17}
\]

由此得：

\[
\mathbf {N} (\operatorname{div} _ {\mathbf {B}} (b)) = \operatorname{div} _ {\mathbf {A}} (\mathbf {N} _ {\mathbf {L} / \mathbf {K}} (b)).\tag{18}
\]

公式 (15) 与 (18) 表明，通过取商，同态 N 与 i 定义了一些同态，对它们我们将同样滥用记号，记作：

\[
\mathrm{N}: \mathrm{C} (\mathrm{B}) \rightarrow \mathrm{C} (\mathrm{A}), \quad i: \mathrm{C} (\mathrm{A}) \rightarrow \mathrm{C} (\mathrm{B}).
\]

注意，同态 \(i\colon\mathbf{C}(\mathbf{A})\to\mathbf{C}(\mathbf{B})\) 一般不是单射（§ 3，习题 7）。

回忆：对每个 B-模 R，\(R_{[A]}\) 表示把 R 的纯量限制到 A 而得到的 A-模（《代数》第 II 章，§1，第 13 节）。

命题 18. (i) R 是伪零 B-模，当且仅当 A-模 \(\mathbf{R}_{[\mathrm{A}]}\) 是伪零的。

\begin{enumerate}
\def\labelenumi{(\roman{enumi})}
\setcounter{enumi}{1}
\tightlist
\item
  R 是有限生成扭 B-模，当且仅当 \(R_{[A]}\) 是有限生成扭 A-模，且此时：
\end{enumerate}

\[
\chi_ {\mathbf {A}} (\mathbf {R} _ {[ \mathbf {A} ]}) = \mathbf {N} (\chi_ {\mathbf {B}} (\mathbf {R})).\tag{19}
\]

\begin{enumerate}
\def\labelenumi{(\roman{enumi})}
\setcounter{enumi}{2}
\tightlist
\item
  R 是有限生成 B-模，当且仅当 \(R_{[A]}\) 是有限生成 A-模，且此时：
\end{enumerate}

\begin{enumerate}
\def\labelenumi{(\arabic{enumi})}
\setcounter{enumi}{19}
\tightlist
\item
\end{enumerate}

\[
c _ {\mathbf {A}} (\mathrm{R} _ {[ \mathbf {A} ]}) = \mathrm{N} (c _ {\mathbf {B}} (\mathrm{R})) + r _ {\mathbf {B}} (\mathrm{R}) c _ {\mathbf {A}} (\mathrm{B})\tag{21}
\]

\[
r _ {\mathrm{A}} (\mathrm{R} _ {[ \mathrm{A} ]}) = n. r _ {\mathrm{B}} (\mathrm{R}) \quad (\text { recall   that } n = r _ {\mathrm{A}} (\mathrm{B})).
\]

由于 B 是有限生成 A-模，R 是有限生成 B-模当且仅当 \(R_{[A]}\) 是有限生成 A-模。此外，若 b 是 R 的零化子，则 \(b \cap A = a\) 是 \(R_{[A]}\) 的零化子；由于 B 在 A 上是整的，除 0 外没有理想位于 A 的理想 0 之上（第 V 章，§2，第 1 节，命题 1 的推论 1），故说 \(\mathfrak{a} \neq 0\) 与说 \(\mathfrak{b} \neq 0\) 是一回事。

\begin{enumerate}
\def\labelenumi{(\roman{enumi})}
\item
  依据这最后的注记，我们可以只考虑 R 是扭 B-模的情形。若 b 含于某个素理想 \(\mathfrak{P} \in \mathrm{P}(B)\) ，则 a 含于 \(\mathfrak{P} \cap A = p\) ，后者高为 1。反之，若 a 含于某个素理想 \(p \in \mathrm{P}(A)\) ，则存在 B 的素理想 P，它包含 b 且位于 p 之上（第 V 章，§2，第 1 节，定理 1 的推论 2）。断言 (i) 即由这些注记及第 4 节定义 2 得出。
\item
  对每个有限生成扭 B-模 R，我们写
\end{enumerate}

\[
\phi (\mathbf {R}) = \chi_ {\mathbf {A}} (\mathbf {R} _ {[ \mathbf {A} ]});
\]

显然（对有限生成扭 B-模）\(\phi\) 满足第 5 节命题 11 的条件 (1) 与 (2)（顾及 (i)）。因此存在同态 \(\theta: \mathrm{D}(\mathrm{B}) \to \mathrm{D}(\mathrm{A})\) ，使对每个有限生成扭 B-模 R 有 \(\phi(\mathrm{R}) = \theta(\chi_{\mathbf{B}}(\mathrm{R}))\) 。同态 \(\theta\) 由它在每个形如 \(\mathrm{B}/\mathfrak{P}\) （ \(\mathfrak{P} \in \mathrm{P}(\mathrm{B})\) ）的 B-模上的值确定，因为 \(\chi_{\mathbf{B}}(\mathrm{B}/\mathfrak{P}) = \mathfrak{P}\) 。现在，对每个素理想 \(\mathfrak{q} \neq \mathfrak{p} = \mathfrak{P} \cap \mathrm{A}\) （ \(\mathrm{P}(\mathrm{A})\) 中），有 \(\mathfrak{p} \notin \mathfrak{q}\) ，从而 \((\mathrm{B}/\mathfrak{P})_{\mathfrak{q}} = 0\) 。另一方面，若置 \(\mathrm{S} = \mathrm{A} - \mathfrak{p}\) ，则 \(\mathfrak{P} \cdot \mathrm{S}^{-1}\mathrm{B}\) 是 \(\mathrm{S}^{-1}\mathrm{B}\) 的极大理想，且 \((\mathrm{B}/\mathfrak{P})_{\mathfrak{p}} = \mathrm{S}^{-1}\mathrm{B}/\mathfrak{P}\) 。\(\mathrm{S}^{-1}\mathrm{B}\) 同构于 \(\mathrm{B}/\mathfrak{P}\) 的分式域（第 II 章，§ 2，第 5 节，命题 11），即赋值 \(v_{\mathfrak{P}}\) 的剩余域；因此它作为 \(\mathrm{A}_{\mathfrak{p}}\) -模的长度是 \(f_{\mathfrak{P}/\mathfrak{p}}\) ；这就证明了 \(\theta = \mathrm{N}\) （第 5 节，定义 4）。

\begin{enumerate}
\def\labelenumi{(\roman{enumi})}
\setcounter{enumi}{2}
\tightlist
\item
  若 T 是 R 的扭子模，则 \(T_{[A]}\) 是 \(R_{[A]}\) 的扭子模，且 \((\mathrm{R}/\mathrm{T})_{[\mathrm{A}]} = \mathrm{R}_{[\mathrm{A}]}/\mathrm{T}_{[\mathrm{A}]}\) ；因此为证 (21)，我们可以只考虑 R 无扭的情形。这时 R 等同于 \(\mathrm{R}_{(\mathrm{L})}\) 的一个子 B-模，且含有一组基 \((e_{i})_{1 \leqslant i \leqslant m}\) —— \(\mathrm{R}_{(\mathrm{L})}\) 在 L 上的基。若 \((b_{j})_{1 \leqslant j \leqslant n}\) 是 L 在 K 上的一组由 B 的元素组成的基，则诸 \(b_{j}e_{i}\) 构成 \(\mathrm{R}_{(\mathrm{L})}\) 在 K 上的一组由 R 的元素组成的基，由此得 (21)。另一方面，设 M 是 R 的自由子 B-模，使 R/M 是扭 B-模；由于 \(M_{[A]}\) 是 \(r_{\mathrm{B}}(\mathrm{R})\) 个同构于 B 的 A-模的直和，由命题 16 (i)
\end{enumerate}

\[
c _ {\mathbf {A}} (\mathbf {M} _ {[ \mathbf {A} ]}) = r _ {\mathbf {B}} (\mathbf {R}). c _ {\mathbf {A}} (\mathbf {B}).
\]

此外，依据 (19)，\(c_{\mathbf{A}}((\mathbf{R} / \mathbf{M})_{[\mathbf{A}]})=-c_{\mathbf{A}}(\mathbf{N}(\chi_{\mathbf{B}}(\mathbf{R} / \mathbf{M})))\) ；但由同态 \(\mathrm{N}\colon\mathrm{C}(\mathrm{B})\to\mathrm{C}(\mathrm{A})\) 的定义，\(c_{\mathbf{A}}(\mathrm{N}(d))=\mathrm{N}(c_{\mathbf{B}}(d))\) 对一切 \(d\in\mathrm{D}(\mathrm{B})\) 成立，且按定义有 \(c_{\mathbf{B}}(\chi(\mathrm{R}/\mathrm{M}))=-c_{\mathbf{B}}(\mathrm{R})\) ，最终得

\[
c _ {\mathbf {A}} ((\mathbf {R} / \mathbf {M}) _ {[ \mathbf {A} ]}) = \mathbf {N} (c _ {\mathbf {B}} (\mathbf {R}));
\]

于是只需应用命题 16 (i) 便得 (20)。

命题 19. 设 R 是有限生成 B-模。R 是自反的，当且仅当 \(R_{[A]}\) 是自反 A-模。

我们在命题 18 的证明中已指出，R 是无扭 B-模当且仅当 \(R_{[A]}\) 是无扭 A-模。因此可以设 R 是 \(W = R \otimes_{B} L\) 关于 B 的格。我们将使用下述引理：

引理 4. 设 W 是 L 上有限秩的向量空间，R 是 W 关于 B 的格。则对一切 \(\mathfrak{p} \in \mathrm{P}(\mathbf{A})\) ，\((\mathbf{R}_{[\mathbf{A}]})_{\mathfrak{p}} = \bigcap_{\mathfrak{P}/\mathfrak{p}} \mathbf{R}_{\mathfrak{P}}\) 。

若 \(S = A - p\) ，则环 \(S^{-1}B\) 的素理想由 \(B\) 中与 \(S\) 不相交的素理想生成，换言之即诸理想 \(\mathfrak{P}_i\) （ \(1 \leqslant i \leqslant m\) ，位于 \(p\) 之上）与理想 (0)；这表明 \(S^{-1}B\) 是半局部环，其极大理想是诸 \(m_i = \mathfrak{P}_i(S^{-1}B)\) （ \(1 \leqslant i \leqslant m\) ）；此外，局部环 \((S^{-1}B)_{\mathfrak{P}_i}\) 同构于 \(B_{m_i}\) （第 II 章，§ 2，第 5 节，命题 11），从而是离散赋值环。于是环 \(S^{-1}B\) 是 Dedekind 整环（§ 2，第 2 节，定理 1 (f)），且由于它是半局部的，它是主理想整环（§ 2，第 2 节，命题 1）。这样，\((R_{[A]})_p\) 就等于 \(S^{-1}R\) （视为 \(A_p\) -模） ；由上所述，\(S^{-1}R\) 是 \(W\) 关于 \(S^{-1}B\) 的自由格，故可对它应用第 2 节的定理 2，得 \(S^{-1}R = \bigcap_i (S^{-1}R)_{m_i}\) ；但 \((S^{-1}R)_{m_i} = R_{\mathfrak{P}_i}\) ，这就证明了引理。

回到命题 19 的证明，由引理 4，

\[
\bigcap_ {\mathfrak {P} \in P (B)} R _ {\mathfrak {P}} = \bigcap_ {\mathfrak {p} \in P (A)} \left(R _ {[ A ]}\right) _ {\mathfrak {p}}
\]

结论由第 2 节的定理 2 得出。

推论. 环 B 是自反 A-模。

命题 20. (i) 有限生成 A-模 M 是伪零的，当且仅当 \(\mathbf{M} \otimes_{\mathbf{A}} \mathbf{B}\) 是伪零 B-模。

\begin{enumerate}
\def\labelenumi{(\roman{enumi})}
\setcounter{enumi}{1}
\tightlist
\item
  若 \(\mathbf{M}\) 是有限生成扭 A-模，则 \(\mathbf{M} \otimes_{\mathbf{A}} \mathbf{B}\) 是有限生成 B-模，且：
\end{enumerate}

\[
\chi_ {\mathbf {B}} (\mathbf {M} \otimes_ {\mathbf {A}} \mathbf {B}) = i (\chi_ {\mathbf {A}} (\mathbf {M})).\tag{22}
\]

\begin{enumerate}
\def\labelenumi{(\roman{enumi})}
\setcounter{enumi}{2}
\tightlist
\item
  若 \(\mathbf{M}\) 是有限生成 A-模，则 \(\mathbf{M} \otimes_{\mathbf{A}} \mathbf{B}\) 是有限生成 B-模，且：
\end{enumerate}

\begin{enumerate}
\def\labelenumi{(\arabic{enumi})}
\setcounter{enumi}{22}
\tightlist
\item
\end{enumerate}

\[
c _ {\mathbf {B}} (\mathbf {M} \otimes_ {\mathbf {A}} \mathbf {B}) = i (c _ {\mathbf {A}} (\mathbf {M}))\tag{24}
\]

\[
r _ {\mathbf {B}} (\mathbf {M} \otimes_ {\mathbf {A}} \mathbf {B}) = r _ {\mathbf {A}} (\mathbf {M}).
\]

\begin{enumerate}
\def\labelenumi{(\roman{enumi})}
\tightlist
\item
  设 \(\mathfrak{P}\) 是 \(\mathbf{B},\mathfrak{p} = \mathfrak{P}\cap \mathbf{A}\) ；则 \((\mathbf{M}\otimes_{\mathbf{A}}\mathbf{B})_{\mathfrak{P}} = \mathbf{M}\otimes_{\mathbf{A}}\mathbf{B}_{\mathfrak{P}}\) （第 II 章，§ 2，第 7 节，命题 18），且另一方面
\end{enumerate}

\[
\mathbf {M} \otimes_ {\mathbf {A}} \mathbf {B} _ {\mathfrak {p}} = (\mathbf {M} \otimes_ {\mathbf {A}} \mathbf {A} _ {p}) \otimes_ {\mathbf {A} _ {p}} \mathbf {B} _ {\mathfrak {p}} = \mathbf {M} _ {p} \otimes_ {\mathbf {A} _ {p}} \mathbf {B} _ {\mathfrak {p}};
\]

因此关系 \(\mathbf{M}_{\mathfrak{p}} = 0\) 等价于 \((\mathbf{M} \otimes_{\mathbf{A}} \mathbf{B})_{\mathfrak{P}} = 0\) （第 II 章，§4，第 4 节，引理 4）。只需把这个注记应用于理想 \(\mathfrak{P} = (0)\) 以及诸理想 \(\mathfrak{P} \in \mathbf{P}(\mathbf{B})\) ，即可证明 (i)，其中用到第 4 节定义 2。

为证 (ii)，我们将使用下述引理：

引理 5. 设 \(\mathbf{M}_1, \mathbf{M}_2\) 是两个有限生成 A-模，\(f: \mathbf{M}_1 \to \mathbf{M}_2\) 是单同态。则 \(f \otimes 1_{\mathbf{B}}: \mathbf{M}_1 \otimes_{\mathbf{A}} \mathbf{B} \to \mathbf{M}_2 \otimes_{\mathbf{A}} \mathbf{B}\) 的核是伪零的。

设 \(\mathfrak{p}\) 是 A 的高 \(\leqslant 1\) 的素理想。则 \((\mathbf{M}_i\otimes_{\mathbf{A}}\mathbf{B})_{\mathfrak{p}} = (\mathbf{M}_i)_{\mathfrak{p}}\otimes_{\mathbf{A}_{\mathfrak{p}}}\mathbf{B}_{\mathfrak{p}}\) \((i = 1,2)\) （第 II 章，§2，第 7 节，命题 18），且 \((f\otimes 1_{\mathbf{B}})_{\mathfrak{p}} = f_{\mathfrak{p}}\otimes 1_{\mathbf{B}_{\mathfrak{p}}}\) ；\(f\) 是单射这一假设蕴含 \(f_{\mathfrak{p}}\) 也是单射（第 II 章，§2，第 4 节，定理 1）；另一方面，鉴于 \(\mathfrak{p},\mathbf{A}_{\mathfrak{p}}\) 是主理想整环，\(\mathbf{B}_{\mathfrak{p}}\) 是有限生成的无扭 \(\mathbf{A}_{\mathfrak{p}}\) -模，从而是自由的；我们断定 \(f_{\mathfrak{p}}\otimes 1_{\mathbf{B}_{\mathfrak{p}}}\) 本身是单射。若 \(\mathbf{I} = \operatorname {Ker}(f\otimes 1)\) ，则 \(\mathbf{I}_{\mathfrak{p}} = \operatorname {Ker}((f\otimes 1)_{\mathfrak{p}})\) （第 II 章，§2，第 4 节，定理 1）；因此 \(\mathrm{I_p} = 0\) ，故更有 \(\mathrm{I}_{\mathfrak{P}} = (\mathrm{I}_{\mathfrak{p}})_{\mathfrak{P}} = 0\) （对 \(\mathfrak{P}| \mathfrak{p}\) ），这就证明了引理（第 4 节，定义 2）。

现在回到 (ii) 的证明。对每个有限生成扭 A-模 M，写 \(\phi(M) = \chi_{B}(M \otimes_{A} B)\) ；由 (i) 可知，若 M 是伪零的，则 \(\phi(M) = 0\) 。另一方面，考虑有限生成扭 A-模的一个正合列：

\[
0 \rightarrow \mathrm{M} _ {1} \rightarrow \mathrm{M} _ {2} \rightarrow \mathrm{M} _ {3} \rightarrow 0.
\]

由引理 4 可知，存在 B-模的一个正合列：

\[
0 \rightarrow \mathrm{I} \rightarrow \mathrm{M} _ {1} \otimes_ {\mathrm{A}} \mathrm{B} \rightarrow \mathrm{M} _ {2} \otimes_ {\mathrm{A}} \mathrm{B} \rightarrow \mathrm{M} _ {3} \otimes_ {\mathrm{A}} \mathrm{B} \rightarrow 0
\]

其中 I 是伪零的。利用第 5 节命题 10 的推论，我们得到 \(\phi(\mathbf{M}_2) = \phi(\mathbf{M}_1) + \phi(\mathbf{M}_3)\) 。于是我们由第 5 节的命题 11 断定：存在同态 \(\theta: \mathbf{D}(\mathbf{A}) \to \mathbf{D}(\mathbf{B})\) ，使对每个有限生成扭 A-模 M 有 \(\phi(\mathbf{M}) = \theta(\chi_{\mathbf{A}}(\mathbf{M}))\) 。为证 \(\theta = i\) ，只需证明 \(\phi(\mathbf{A}/\mathfrak{p}) = i(\mathfrak{p})\) 对一切 \(\mathfrak{p} \in \mathbf{P}(\mathbf{A})\) 成立；现在 \((\mathbf{A}/\mathfrak{p}) \otimes_{\mathbf{A}} \mathbf{B} = \mathbf{B}/\mathfrak{p}\mathbf{B}\) ，且对一切 \(\mathfrak{P} \in \mathbf{P}(\mathbf{B})\) ，\((\mathbf{B}/\mathfrak{p}\mathbf{B})_{\mathfrak{P}} = \mathbf{B}_{\mathfrak{P}}/\mathfrak{p}\mathbf{B}_{\mathfrak{P}}\) ；若 \(\mathfrak{P}\) 不位于 \(\mathfrak{p}\) 之上，则最后的模是 0 ；反之，若 \(\mathfrak{P}| \mathfrak{p}, \mathbf{B}_{\mathfrak{P}}/\mathfrak{p}\mathbf{B}_{\mathfrak{P}}\) 是 \(\mathbf{B}_{\mathfrak{P}}\) -模，其长度为 \(e(\mathfrak{P}/\mathfrak{p})\) （按分歧指数的定义，第 VI 章，§ 8，第 1 节）；

因此 \(\chi_{\mathbf{B}}(\mathrm{B} / p\mathrm{B}) = \sum_{\mathfrak{P} / p}e_{\mathfrak{P} / p}\cdot \mathfrak{P} = i(p)\) ，这就完成了 (ii) 的证明。

公式 (24) 是直接的，因为

\[
(\mathbf {M} \otimes_ {\mathbf {A}} \mathbf {B}) \otimes_ {\mathbf {B}} \mathbf {L} = \mathbf {M} \otimes_ {\mathbf {A}} \mathbf {L} = (\mathbf {M} \otimes_ {\mathbf {A}} \mathbf {K}) \otimes_ {\mathbf {K}} \mathbf {L}
\]

且 \((\mathbf{M} \otimes_{\mathbf{A}} \mathbf{K}) \otimes_{\mathbf{K}} \mathbf{L}\) 在 L 上的秩等于 \(M \otimes_{A} K\) 在 K 上的秩。为证 (23)，考虑 M 的自由子模 H，使 Q = M/H 是扭 A-模。如上应用引理 4，我们得到 B-模的一个正合列：

\[
0 \rightarrow \mathrm{I} \rightarrow \mathrm{H} \otimes_ {\mathrm{A}} \mathrm{B} \rightarrow \mathrm{M} \otimes_ {\mathrm{A}} \mathrm{B} \rightarrow \mathrm{Q} \otimes_ {\mathrm{A}} \mathrm{B} \rightarrow 0
\]

因此由第 7 节命题 16 (ii)、(v) 以及命题 16 的推论 1 得

\[
c _ {\mathrm{B}} (\mathrm{M} \otimes_ {\mathrm{A}} \mathrm{B}) = c _ {\mathrm{B}} (\mathrm{Q} \otimes_ {\mathrm{A}} \mathrm{B}) = - c _ {\mathrm{B}} \left(\chi_ {\mathrm{B}} (\mathrm{Q} \otimes_ {\mathrm{A}} \mathrm{B})\right) = - c _ {\mathrm{B}} (i \left(\chi_ {\mathrm{A}} (\mathrm{Q})\right))
\]

这是依据 (ii)；但由同态 \(i: \mathbf{C}(\mathbf{A}) \to \mathbf{C}(\mathbf{B})\) 的定义，\(c_{\mathbf{B}}(i(\chi_{\mathbf{A}}(\mathbf{Q}))) = i(c_{\mathbf{A}}(\chi_{\mathbf{A}}(\mathbf{Q}))) = -i(c_{\mathbf{A}}(\mathbf{M}))\) ，这就完成了 (23) 的证明。

\textbf{注}\par

\begin{enumerate}
\def\labelenumi{(\arabic{enumi})}
\item
  若 M 是自反 A-模，则 \(M \otimes_{A} B\) 未必自反（习题 6）。但当 B 是平坦 A-模时它是自反的（第 2 节，命题 8）。
\item
  设 C 是第三个整闭 Noether 整环，满足 \(B \subset C\) ，且 C 是有限生成 B-模（从而也是有限生成 A-模）。则我们有传递公式：
\item
\end{enumerate}

\[
\mathbf {N} _ {\mathbf {C} / \mathbf {A}} = \mathbf {N} _ {\mathbf {B} / \mathbf {A}} \circ \mathbf {N} _ {\mathbf {C} / \mathbf {B}},\tag{26}
\]

\[
i _ {\mathbf {C} / \mathbf {A}} = i _ {\mathbf {C} / \mathbf {B}} \circ i _ {\mathbf {B} / \mathbf {A}}
\]

它们由分歧指数与剩余类次数的传递公式（第 VI 章，§ 8，第 1 节，引理 1）立即得出。

